\chapter{Réalisation}

\section{Fonctionnalités réalisées (Lot A)}
\begin{itemize}
  \item Implémentation complète du backend Spring Boot (controllers, services, repositories).
  \item Développement des modules administratifs : demandes, absences, justifications, réclamations, notifications.
  \item Intégration de la persistance avec JPA/Hibernate et gestion des profils de base de données.
  \item Mise en place de la gestion d'erreurs globale et standardisation des réponses API.
  \item Validation fonctionnelle de l'API via tests manuels (Postman) et vérification des scénarios critiques.
\end{itemize}

\section{Résultats techniques obtenus}
Le lot backend a permis d'obtenir une API exploitable par le frontend avec un nombre élevé d'endpoints (environ 90+), couvrant les fonctionnalités principales du système. Les services sont organisés selon une logique métier claire et une architecture maintenable.

Les points notables sont :
\begin{itemize}
  \item hiérarchie d'utilisateurs (admin/enseignant/étudiant) ;
  \item gestion des départements, cours, matières, absences et demandes ;
  \item workflows administratifs suivis de bout en bout ;
  \item base de code prête à l'évolution (sécurité, tests automatiques, monitoring).
\end{itemize}

\section{Contraintes rencontrées et solutions apportées}
\begin{tabularx}{\textwidth}{lX}
\toprule
\textbf{Contrainte} & \textbf{Solution appliquée} \\
\midrule
Évolution du modèle de données & Refactoring des entités et consolidation des relations JPA pour éviter les incohérences. \\
Hétérogénéité des réponses API & Normalisation des retours via un format de réponse commun et gestion centralisée des erreurs. \\
Synchronisation frontend/backend & Stabilisation des contrats d'API et validation systématique des endpoints critiques. \\
Qualité du code & Structuration stricte en couches (Controller/Service/Repository) pour améliorer la maintenabilité. \\
\bottomrule
\end{tabularx}

\section{Compétences acquises}
Au cours de ce lot, les compétences suivantes ont été consolidées :
\begin{itemize}
  \item conception d'API REST avec Spring Boot ;
  \item modélisation relationnelle et mapping JPA/Hibernate ;
  \item organisation d'un backend en architecture multicouche ;
  \item gestion des workflows administratifs ;
  \item diagnostic et résolution de problèmes d'intégration applicative.
\end{itemize}

\section{Livrables du Lot A}
\begin{itemize}
  \item code source backend (modules model/repository/service/controller) ;
  \item endpoints API documentés ;
  \item jeux de tests manuels Postman ;
  \item structure de données cohérente pour intégration frontend.
\end{itemize}

\section{Division des tâches du projet}
\begin{tabularx}{\textwidth}{lX}
\toprule
\textbf{Composant} & \textbf{Responsabilité Lot A (ce rapport)} \\
\midrule
Backend & Architecture et logique métier Java/Spring. \\
Base de données & Modélisation, mapping JPA et cohérence des données. \\
API & Endpoints REST, validation et gestion des exceptions. \\
Processus Admin & Implémentation complète des workflows administratifs. \\
Intégration & Mise à disposition des services consommés par le frontend (contrats API stables). \\
\bottomrule
\end{tabularx}
